\documentclass[../main.tex]{subfiles}\begin{document}
\section{Risikoanalyse}

For at sikre projektets fremdrift udføres en risikoanalyse, der kortlægger de mest kritiske forhindringer. Målet er at etablere klare strategier for risikohåndtering, hvilket kan indebære alt fra direkte problemløsning til en justering af projektets samlede omfang.\newline

Analysen systematiseres i en risikomatrix, hvor de enkelte risici tildeles en "risk score". Denne vurdering er resultatet af en direkte afvejning mellem sandsynligheden for, at udfordringen opstår, og konsekvensens alvor, såfremt risikoen realiseres
%TC:ignore
\begin{table}[H]
    \caption{Risiko matrix}
    \label{tab:risikomatrirex}
    \centering
    \renewcommand{\arraystretch}{1.5}
    \begin{tabular}{|c|c|c|c|c|c|}
    
    \hline
    & \multicolumn{5}{c|}{\textbf{Impact}} \\ \cline{2-6}
    \textbf{Likelihood} & \textbf{Negligible} & \textbf{Minor} & \textbf{Moderate} & \textbf{Significant} & \textbf{Severe} \\ \hline
    \textbf{Very likely}  & \cellcolor{green!50}Low Med & \cellcolor{yellow!30}Medium & \cellcolor{orange!50}Med Hi & \cellcolor{red!60}High & \cellcolor{red!60}High \\ \hline
    \textbf{Likely}       & \cellcolor{green!30}Low & \cellcolor{green!50}Low Med & \cellcolor{yellow!30}Medium & \cellcolor{orange!50}Med Hi & \cellcolor{red!60}High \\ \hline
    \textbf{Possible}     & \cellcolor{green!30}Low & \cellcolor{green!50}Low Med & \cellcolor{yellow!30}Medium & \cellcolor{orange!50}Med Hi & \cellcolor{orange!50}Med Hi \\ \hline
    \textbf{Unlikely}     & \cellcolor{green!30}Low & \cellcolor{green!50}Low Med & \cellcolor{yellow!30}Medium & \cellcolor{yellow!30}Medium & \cellcolor{orange!50}Med Hi \\ \hline
    \textbf{Very Unlikely}& \cellcolor{green!30}Low & \cellcolor{green!30}Low & \cellcolor{green!50}Low Med & \cellcolor{yellow!30}Medium & \cellcolor{yellow!30}Medium \\ \hline
    
    \end{tabular}
\end{table}
%TC:endignore
\newpage
\subsection*{Processmæssige forhold}
%TC:ignore
\begin{longtable}{| P{3.5cm} | P{2cm} | P{2cm} | P{2.5cm} | P{3cm} |}
    \caption{Risiko tabel - Processmæssige forhold} \\
    \hline
    \rowcolor{gray!50}
    \textbf{Risk} & \textbf{Likelihood} & \textbf{Impact} & \textbf{Risk Rating} & \textbf{Response} \\
    \hline
    \endfirsthead % Dette markerer slutningen af den første header

    \multicolumn{5}{c}%
    {{\bfseries Tabel \thetable{} -- fortsat fra forrige side}} \\
    \hline
    \rowcolor{gray!50}
    \textbf{Risk} & \textbf{Likelihood} & \textbf{Impact} & \textbf{Risk Rating} & \textbf{Response} \\
    \hline
    \endhead % Dette markerer slutningen af headeren for de efterfølgende sider

    \hline \multicolumn{5}{|r|}{{Fortsat på næste side}} \\ \hline
    \endfoot % Dette markerer slutningen af foden for de fleste sider

    \hline
    \endlastfoot % Dette markerer slutningen af foden for den sidste side

    Risiko 1: Frafald fra gruppen & Very Unlikely & Significant & \cellcolor{yellow!30}Medium & Gruppen holder løbende standup møder for at vidensdele, samt plejer relationen uden for projektet. \\
    \hline
    \rowcolor{gray!30}
    Risiko 2: Samarbejdsudfordringer & Possible & Significant & \cellcolor{orange!50}Med Hi & Der er udarbejdet en klar samarbejdskontrakt inden projektets start, og i led af gruppens størrelse, vil udfordringer hurtigt opdages og behandles. \\
    \hline
    Risiko 3: For ambitiøse mål & Possible & Moderate & \cellcolor{yellow!30}Medium & Gruppen justerer mål, så projektet kan færdiggøres i tide på tilfredsstillende vis. \\
    \hline
    \rowcolor{gray!30}
    Risiko 4: Manglende kongruens i projektets mål blandt gruppens medlemmer & Possible & Moderate & \cellcolor{yellow!30}Medium & Gruppen holder møde om, hvad projektet skal opnå, og på denne måde får vi genafstemt hvilke mål vi vil opnå. \\
    \hline
\end{longtable}
%TC:endignore
\newpage
\subsection*{Teknologiske forhold}
%TC:ignore

\begin{longtable}{| P{3.5cm} | P{2cm} | P{2cm} | P{2.5cm} | P{3cm} |}
    \caption{Risiko tabel - teknologiske forhold} \\
    \hline
    \rowcolor{gray!50}
    \textbf{Risk} & \textbf{Likelihood} & \textbf{Impact} & \textbf{Risk Rating} & \textbf{Response} \\
    \hline
    \endfirsthead % Dette markerer slutningen af den første header

    \multicolumn{5}{c}%
    {{\bfseries Tabel \thetable{} -- fortsat fra forrige side}} \\
    \hline
    \rowcolor{gray!50}
    \textbf{Risk} & \textbf{Likelihood} & \textbf{Impact} & \textbf{Risk Rating} & \textbf{Response} \\
    \hline
    \endhead % Dette markerer slutningen af headeren for de efterfølgende sider

    \hline \multicolumn{5}{|r|}{{Fortsat på næste side}} \\ \hline
    \endfoot % Dette markerer slutningen af foden for de fleste sider

    \hline
    \endlastfoot % Dette markerer slutningen af foden for den sidste side

    Risiko 1: Datakonflikter i grænsefladen mellem frontend og backend & Possible & Moderate & \cellcolor{yellow!30}Medium & API-endepunkter defineres via en fælles kontrakt (f.eks. OpenAPI/Swagger) inden implementering. Der anvendes DTO'er for at sikre en konsekvent datastruktur. \\
    \hline
    \rowcolor{gray!30}
    Risiko 2: Kritiske teknologiske begrænsninger opdages sent & Unlikely & Severe & \cellcolor{orange!50}Med Hi & Inden den endelige teknologiske arkitektur fastlægges, udarbejdes en gennemgående teknisk analyse og et tidligt Proof of Concept (PoC) på projektets mest komplekse funktion. \\
    \hline
    Risiko 3: Visuelle rendering-fejl på tværs af forskellige browsere og skærmstørrelser & Possible & Minor & \cellcolor{green!50}Low Med & Der anvendes et responsivt CSS-framework. Mindre kosmetiske afvigelser accepteres systematisk (Risk Acceptance), så længe systemets kernefunktionalitet er intakt. \\
    \hline
    \rowcolor{gray!30}
    Risiko 4: Udfordringer med opsætning af cloud/deployment-miljø & Likely & Moderate & \cellcolor{yellow!30}Medium & Deployment-pipelinen opsættes i starten af projektet (eks. via GitHub Actions), så koden løbende integreres og testes \\
    \hline
    Risiko 5: Skemakonflikter mellem applikationens domænemodel og databasen & Possible & Moderate & \cellcolor{yellow!30}Medium & Databaseændringer styres gennem automatiserede migrations-scripts (f.eks. EF Core Migrations), og dataadgangen valideres via integrationstests. \\
    \hline
    \rowcolor{gray!30}
    Risiko 6: Sikkerhedsbrister i brugerhåndtering (Autentificering/Autorisation) & Unlikely & Significant & \cellcolor{yellow!30}Medium & I stedet for at bygge et system fra bunden, integreres en anerkendt standardløsning (f.eks. OAuth2 eller et etableret bibliotek). \\
    \hline
    Risiko 7: Inkonsistens i data ved samtidige forespørgsler (Concurrency) & Possible & moderate & \cellcolor{yellow!30}Medium & Databasen designes med eksplicitte transaktioner og korrekte isolationsniveauer. \\
    \hline
\end{longtable}
%TC:endignore
\newpage

\subsection*{Risici Konklusion}

Eftersom ingen af de identificerede risici overstiger niveauet "Medium High", vurderes det overordnede risikobillede at være acceptabelt. Projektet kan derfor fortsætte inden for de overordnede rammer, dog med den betingelse, at de definerede mitigeringsstrategier integreres aktivt i arbejdsgangen. Risikoanalysen viser, at gruppens primære opmærksomhedspunkt er udvælgelsen og valideringen af teknologistakken.

\end{document}
